\section{Fijaciones consecutivas en misma celda}
\label{sec:fijaciones_misma_celda}

Como consecuencia de la diferencia de resolución espacial entre el 
modelo nnELM, que opera sobre una grilla de celdas, y el registro 
de eye-tracking, que captura fijaciones a nivel de píxel, puede 
ocurrir que dos fijaciones consecutivas de un mismo trial sean 
asignadas a la misma celda de la grilla. En ese caso, surge la 
pregunta de si ambas fijaciones recibirían valores idénticos para 
todas las features, lo que introduciría repetición artificial en 
los predictores del modelo de regresión. Para evaluar la magnitud 
y el impacto de este fenómeno, se realizó un análisis específico 
de los pares de fijaciones consecutivas que caen dentro de la 
misma celda.

Del total de 
$61.252$ fijaciones, $2.908$ ($4.7\%$) comparten celda con la 
fijación inmediatamente anterior. Este porcentaje varía entre 
sujetos (rango: $0.6\%$--$9.6\%$), sin que ningún participante 
presente una proporción que comprometa el análisis.

La Tabla~\ref{tab:misma_celda} presenta, para cada feature, el número 
de pares con valores idénticos y distintos dentro de los $2.908$ pares 
de fijaciones consecutivas en la misma celda.

\begin{table}[H]
\centering
\caption{Verificación de identidad de valores en pares de fijaciones 
consecutivas que caen en la misma celda ($N = 2.908$ pares).}
\label{tab:misma_celda}
\small
\setlength{\tabcolsep}{8pt}
\renewcommand{\arraystretch}{1.2}
\begin{tabular}{lrrr}
\toprule
\textbf{Feature} & \textbf{Idénticas} & \textbf{Distintas} & \textbf{\% idénticas} \\
\midrule
\texttt{saliencia\_media}               & 2.908 & 0     & 100.0\% \\
\texttt{saliencia\_mediana}             & 2.908 & 0     & 100.0\% \\
\texttt{similitud\_media}               & 2.908 & 0     & 100.0\% \\
\texttt{similitud\_mediana}             & 2.908 & 0     & 100.0\% \\
\texttt{entropia\_prior\_saliencia}     & 2.908 & 0     & 100.0\% \\
\midrule
\texttt{competencia\_densidad}          & 1.943 & 965   & 66.8\%  \\
\texttt{competencia\_densidad\_norm}    & 1.943 & 965   & 66.8\%  \\
\texttt{saliencia\_pixel}               & 1.596 & 1.312 & 54.9\%  \\
\texttt{similitud\_pixel}               & 1.101 & 1.807 & 37.9\%  \\
\texttt{distancia\_grilla}              & 197   & 2.711 & 6.8\%   \\
\texttt{distancia\_pixels}              & 197   & 2.711 & 6.8\%   \\
\texttt{gap\_entropia}                  & 24    & 2.884 & 0.8\%   \\
\midrule
\texttt{evidencia\_visual}              & 0     & 2.908 & 0.0\%   \\
\texttt{competencia\_fp\_var}           & 0     & 2.908 & 0.0\%   \\
\texttt{competencia\_fp\_media}         & 0     & 2.908 & 0.0\%   \\
\texttt{competencia\_fp\_max}           & 0     & 2.908 & 0.0\%   \\
\texttt{posterior}                      & 0     & 2.908 & 0.0\%   \\
\texttt{entropia\_posterior}            & 0     & 2.908 & 0.0\%   \\
\texttt{ganancia\_informacion}          & 0     & 2.908 & 0.0\%   \\
\texttt{competencia\_ratio}             & 0     & 2.908 & 0.0\%   \\
\bottomrule
\end{tabular}
\end{table}

El primer grupo corresponde a las features con valores idénticos en 
el $100\%$ de los pares: \texttt{saliencia\_media}, 
\texttt{saliencia\_mediana}, \texttt{similitud\_media}, 
\texttt{similitud\_mediana} y \texttt{entropia\_prior\_saliencia}. 
Estas features dependen únicamente de la celda y de la imagen, por 
lo que dos fijaciones en la misma celda reciben necesariamente el 
mismo valor. Este comportamiento es esperable y no representa un 
problema metodológico, sino una consecuencia directa de su naturaleza 
estática.

El segundo grupo corresponde a las features con valores siempre 
distintos ($0\%$ de idénticas): \texttt{evidencia\_visual}, 
\texttt{competencia\_fp\_var}, \texttt{competencia\_fp\_media}, 
\texttt{competencia\_fp\_max}, \texttt{posterior}, 
\texttt{entropia\_posterior}, \texttt{ganancia\_informacion}, 
\texttt{gap\_entropia} y \texttt{competencia\_ratio}. Para estas 
features, los mapas internos del modelo nnELM se actualizan en 
cada fijación independientemente de la celda visitada, por lo que 
dos fijaciones consecutivas en la misma celda producen valores 
distintos.

El tercer grupo corresponde a features con comportamiento intermedio, 
que merecen una explicación particular. \texttt{saliencia\_pixel} 
($54.9\%$) y \texttt{similitud\_pixel} ($37.9\%$) dependen del 
píxel exacto de fijación dentro de la celda: dos fijaciones en la 
misma celda pueden caer en píxeles distintos y recibir valores 
distintos, o en el mismo píxel y recibir valores idénticos. 
\texttt{competencia\_densidad} y \texttt{competencia\_densidad\_norm} 
presentan un $66.8\%$ de valores idénticos no por ser estáticas sino 
por su distribución discreta: dado que la mediana es $1$ y el Q75 
es $2$, es frecuente que dos fijaciones consecutivas tengan el mismo 
valor por coincidencia. Finalmente, \texttt{distancia\_grilla} y 
\texttt{distancia\_pixels} presentan un $6.8\%$ de valores idénticos, 
correspondientes a los casos en que el modelo predice la misma celda 
óptima en dos fijaciones consecutivas.

En conjunto, dado que el fenómeno de fijaciones consecutivas en la 
misma celda afecta al $4.7\%$ del total de fijaciones, y que las 
features dinámicas no presentan valores idénticos en esos pares, 
se optó por conservar todas las fijaciones sin tratamiento adicional.

